\chapter{A Complete Worked Example}
\label{chap:ch22}

\section{Synthesizing the Pipeline with Concrete Numbers}
Throughout this textbook, we have examined remote sensing physics, tensor dimensions, attention mathematics, and loss formulations in isolation.

In this chapter, we assemble every concept into a single, complete \textbf{end-to-end numerical walkthrough}. We will take a tiny $2 \times 2$ pixel image, pass it through a miniature multimodal network, compute the composite loss, derive exact backpropagation gradients using the Chain Rule, and execute a parameter update.

\section{The Micro-World Setup}
To keep the manual arithmetic fully transparent while preserving multimodal properties, we define:
\begin{itemize}
    \item Spatial grid: $H = 2, W = 2$ ($2 \times 2$ pixels).
    \item Modalities:
    \begin{itemize}
        \item Sentinel-1 SAR VV ($\mathbf{X}_{\text{SAR}} \in \mathbb{R}^{1 \times 2 \times 2}$).
        \item Sentinel-2 Optical Red band ($\mathbf{X}_{\text{opt}} \in \mathbb{R}^{1 \times 2 \times 2}$).
        \item Soft Cloud Mask ($\mathbf{M} \in \mathbb{R}^{1 \times 2 \times 2}$).
    \end{itemize}
    \item Target ground truth: $\mathbf{Y}^* \in \mathbb{R}^{1 \times 2 \times 2}$.
\end{itemize}

\subsection{Input Data Values}
\[
\mathbf{X}_{\text{SAR}} = \begin{bmatrix} 0.20 & 0.80 \\ 0.20 & 0.80 \end{bmatrix} \quad (\text{left col: calm water; right col: urban concrete})
\]
\[
\mathbf{X}_{\text{opt}} = \begin{bmatrix} 0.10 & 0.95 \\ 0.10 & 0.95 \end{bmatrix} \quad (\text{left col: clear water; right col: obscured by cloud!})
\]
\[
\mathbf{M} = \begin{bmatrix} 0.00 & 1.00 \\ 0.00 & 1.00 \end{bmatrix} \quad (\text{left col: clear sky; right col: thick cloud})
\]
\[
\mathbf{Y}^* = \begin{bmatrix} 0.10 & 0.25 \\ 0.10 & 0.25 \end{bmatrix} \quad (\text{true ground truth: water on left, urban asphalt on right})
\]

\section{The Miniature Multimodal Architecture}
We construct a minimal model with two learnable parameters:
\begin{enumerate}
    \item A radar stem weight $w_{\text{sar}}$ and an optical stem weight $w_{\text{opt}}$:
    \begin{equation}
    \mathbf{P}_{\text{dir}} = w_{\text{sar}} \mathbf{X}_{\text{SAR}} + w_{\text{opt}} \mathbf{X}_{\text{opt}}
    \end{equation}
    \item Initial parameter values: $w_{\text{sar}} = 0.50$, $w_{\text{opt}} = 0.50$.
    \item Blended output:
    \begin{equation}
    \hat{\mathbf{Y}} = (1 - \mathbf{M}) \odot \mathbf{X}_{\text{opt}} + \mathbf{M} \odot \mathbf{P}_{\text{dir}}
    \end{equation}
    \item Objective: $L_1$ reconstruction loss: $\mathcal{L} = \frac{1}{4} \sum_{i=1}^4 |\hat{\mathbf{Y}}_i - \mathbf{Y}^*_i|$.
    \item Learning rate: $\eta = 0.10$.
\end{enumerate}

\section{Step 1: The Forward Pass}
First, compute the network's direct prediction $\mathbf{P}_{\text{dir}}$:
\begin{align*}
\mathbf{P}_{\text{dir}}(0, 0) &= (0.50)(0.20) + (0.50)(0.10) = 0.10 + 0.05 = \mathbf{0.15} \\
\mathbf{P}_{\text{dir}}(0, 1) &= (0.50)(0.80) + (0.50)(0.95) = 0.40 + 0.475 = \mathbf{0.875} \\
\mathbf{P}_{\text{dir}}(1, 0) &= (0.50)(0.20) + (0.50)(0.10) = \mathbf{0.15} \\
\mathbf{P}_{\text{dir}}(1, 1) &= (0.50)(0.80) + (0.50)(0.95) = \mathbf{0.875}
\end{align*}
\[
\mathbf{P}_{\text{dir}} = \begin{bmatrix} 0.150 & 0.875 \\ 0.150 & 0.875 \end{bmatrix}
\]

Next, apply the soft-mask blend gate $\hat{\mathbf{Y}} = (1 - \mathbf{M})\mathbf{X}_{\text{opt}} + \mathbf{M}\mathbf{P}_{\text{dir}}$:
\begin{align*}
\hat{\mathbf{Y}}(0, 0) &= (1 - 0.0)(0.10) + (0.0)(0.15) = \mathbf{0.100} \quad (\text{exact clear ground truth preserved!}) \\
\hat{\mathbf{Y}}(0, 1) &= (1 - 1.0)(0.95) + (1.0)(0.875) = \mathbf{0.875} \quad (\text{cloudy pixel replaced by prediction}) \\
\hat{\mathbf{Y}}(1, 0) &= \mathbf{0.100} \\
\hat{\mathbf{Y}}(1, 1) &= \mathbf{0.875}
\end{align*}
\[
\hat{\mathbf{Y}} = \begin{bmatrix} 0.100 & 0.875 \\ 0.100 & 0.875 \end{bmatrix}
\]

\section{Step 2: Loss Evaluation}
Compare blended prediction $\hat{\mathbf{Y}}$ with ground-truth $\mathbf{Y}^*$:
\[
\hat{\mathbf{Y}} - \mathbf{Y}^* = \begin{bmatrix} 0.100 - 0.100 & 0.875 - 0.250 \\ 0.100 - 0.100 & 0.875 - 0.250 \end{bmatrix} = \begin{bmatrix} 0.000 & +0.625 \\ 0.000 & +0.625 \end{bmatrix}
\]
Evaluating the scalar $L_1$ loss:
\[
\mathcal{L} = \frac{1}{4} \big( |0.000| + |0.625| + |0.000| + |0.625| \big) = \frac{1.250}{4} = \mathbf{0.3125}
\]

\section{Step 3: The Backward Pass (Autograd Derivatives)}
We evaluate the gradient of $\mathcal{L}$ with respect to the blended output pixels:
\[
\frac{\partial \mathcal{L}}{\partial \hat{\mathbf{Y}}} = \frac{1}{4} \text{sign}(\hat{\mathbf{Y}} - \mathbf{Y}^*) = \begin{bmatrix} 0.000 & +0.250 \\ 0.000 & +0.250 \end{bmatrix}
\]
Now, propagate backward through the soft-mask gate:
\[
\frac{\partial \mathcal{L}}{\partial \mathbf{P}_{\text{dir}}} = \frac{\partial \mathcal{L}}{\partial \hat{\mathbf{Y}}} \odot \mathbf{M} = \begin{bmatrix} (0.000)(0.0) & (0.250)(1.0) \\ (0.000)(0.0) & (0.250)(1.0) \end{bmatrix} = \begin{bmatrix} 0.000 & +0.250 \\ 0.000 & +0.250 \end{bmatrix}
\]
Notice that the left column has zero gradient because $\mathbf{M} = 0$!

Now, propagate backward into the weights:
\begin{align*}
\frac{\partial \mathcal{L}}{\partial w_{\text{sar}}} &= \sum_{i,j} \frac{\partial \mathcal{L}}{\partial \mathbf{P}_{\text{dir}}(i,j)} \cdot \mathbf{X}_{\text{SAR}}(i,j) \\
&= (0.000)(0.20) + (0.250)(0.80) + (0.000)(0.20) + (0.250)(0.80) \\
&= 0.000 + 0.200 + 0.000 + 0.200 = \mathbf{+0.400}
\end{align*}
\begin{align*}
\frac{\partial \mathcal{L}}{\partial w_{\text{opt}}} &= \sum_{i,j} \frac{\partial \mathcal{L}}{\partial \mathbf{P}_{\text{dir}}(i,j)} \cdot \mathbf{X}_{\text{opt}}(i,j) \\
&= (0.000)(0.10) + (0.250)(0.95) + (0.000)(0.10) + (0.250)(0.95) \\
&= 0.000 + 0.2375 + 0.000 + 0.2375 = \mathbf{+0.475}
\end{align*}

\section{Step 4: The Optimizer Parameter Update}
Both gradients are positive, meaning the predictions were too bright and the weights must decrease:
\[
w_{\text{sar, new}} = w_{\text{sar}} - \eta \frac{\partial \mathcal{L}}{\partial w_{\text{sar}}} = 0.500 - (0.10)(+0.400) = 0.500 - 0.040 = \mathbf{0.460}
\]
\[
w_{\text{opt, new}} = w_{\text{opt}} - \eta \frac{\partial \mathcal{L}}{\partial w_{\text{opt}}} = 0.500 - (0.10)(+0.475) = 0.500 - 0.0475 = \mathbf{0.4525}
\]

\section{Step 5: Verifying the Update on Step 2}
Let us test the updated weights on a new forward pass:
\[
\mathbf{P}_{\text{dir, new}}(0, 1) = (0.460)(0.80) + (0.4525)(0.95) = 0.368 + 0.4299 = \mathbf{0.7979}
\]
The error on the cloudy pixel dropped from $|0.875 - 0.25| = 0.625$ down to $|0.7979 - 0.25| = \mathbf{0.5479}$!
The new loss is:
\[
\mathcal{L}_{\text{new}} = \frac{1}{4}(0 + 0.5479 + 0 + 0.5479) = \mathbf{0.2739} < 0.3125
\]
The system learned!

\section{Summary of Chapter 22}
\begin{enumerate}
    \item Concrete matrix arithmetic traces the exact pathway from raw numbers through the blend gate to backpropagation.
    \item The mask multiplication $\frac{\partial \mathcal{L}}{\partial \mathbf{P}_{\text{dir}}} = \frac{\partial \mathcal{L}}{\partial \hat{\mathbf{Y}}} \odot \mathbf{M}$ zeroes out upstream gradients on clear pixels.
    \item Gradients accumulate the sensitivity of every cloudy pixel to each internal weight.
    \item Gradient descent strictly reduces reconstruction error on the subsequent step.
\end{enumerate}

\section{Practice Exercises}
\begin{enumerate}
    \item In Step 3, why was the gradient with respect to the clear pixel $(0, 0)$ equal to zero?
    \item If we had used Mean Squared Error ($\mathcal{L} = \frac{1}{2}(\hat{y} - y^*)^2$) instead of $L_1$, what would have been the derivative $\frac{\partial \mathcal{L}}{\partial \hat{y}(0, 1)}$?
    \item If the radar weight received a $2\times$ differential learning rate ($\eta_{\text{sar}} = 0.20$), calculate the updated $w_{\text{sar, new}}$.
\end{enumerate}

\subsection*{Solutions}
\begin{enumerate}
    \item At pixel $(0, 0)$, the mask value was $\mathbf{M} = 0.0$. Multiplying by $\mathbf{M}$ zeroes the gradient flowing to $\mathbf{P}_{\text{dir}}$. Because the original optical pixel already matched the target ($\hat{y} = y^* = 0.10$), the reconstruction error was also zero.
    \item For MSE: $\frac{\partial \mathcal{L}}{\partial \hat{y}} = (\hat{y} - y^*) = (0.875 - 0.250) = \mathbf{+0.625}$.
    \item $w_{\text{sar, new}} = 0.500 - (0.20)(+0.400) = 0.500 - 0.080 = \mathbf{0.420}$.
\end{enumerate}
